\section*{Chapitre \ref{ch:g4:world-physical-map} --- La carte physique du monde}

\begin{solution}{exo:g4:world-physical-map:1}
L'\omterm{def:g4:world-physical-map:altitude}{altitude} d'un lieu est sa hauteur au-dessus du \omterm{def:g4:world-physical-map:altitude}{niveau de la mer}, en
mètres. On la mesure à partir du \omterm{def:g4:world-physical-map:altitude}{niveau de la mer}, le niveau moyen de la
surface de la \omterm{def:g2:water-on-earth:ocean}{mer}.
\end{solution}

\begin{solution}{exo:g4:world-physical-map:2}
(a) Vert foncé~: entre 0 et 200 mètres. (b) Jaune pâle~: entre 500 et
1\,000 mètres. (c) Presque blanc~: plus de 4\,000 mètres.
\end{solution}

\begin{solution}{exo:g4:world-physical-map:3}
(a) L'Amérique du Sud. (b) L'Asie. (c) L'Afrique. (d) L'Amérique du Nord.
\end{solution}

\begin{solution}{exo:g4:world-physical-map:4}
La plus haute \omterm{def:g1:landscapes-around-us:landform}{montagne} est l'Everest (8\,849 mètres). Le plus grand \omterm{def:g1:landscapes-around-us:water}{lac} est
la \omterm{def:g2:water-on-earth:ocean}{mer} Caspienne~; le \omterm{def:g1:landscapes-around-us:water}{lac} le plus profond est le \omterm{def:g1:landscapes-around-us:water}{lac} Baïkal.
\end{solution}

\begin{solution}{exo:g4:world-physical-map:5}
Plus le bleu est foncé, plus la \omterm{def:g2:water-on-earth:ocean}{mer} est profonde~: la \omterm{def:g2:water-on-earth:ocean}{mer} est peu profonde
près des \omterm{def:g1:landscapes-around-us:water}{côtes} (moins de 200 mètres) et bien plus profonde au milieu des
\omterm{def:g2:water-on-earth:ocean}{océans} (plus de 4\,000 mètres en de nombreux endroits).
\end{solution}

\begin{solution}{exo:g4:world-physical-map:6}
La \omterm{def:g4:world-physical-map:physical-map}{carte physique} montre l'\omterm{def:g4:world-physical-map:altitude}{altitude} de la surface. En Antarctique et au
Groenland, la surface est le dessus d'une épaisse couche de glace, bien
au-dessus du \omterm{def:g4:world-physical-map:altitude}{niveau de la mer}~: ils prennent donc les couleurs des terres
hautes.
\end{solution}

\begin{solution}{exo:g4:world-physical-map:7}
Le fond de l'abîme Challenger (près de 11\,000 mètres sous le \omterm{def:g4:world-physical-map:altitude}{niveau de la
mer})~; le fond moyen de l'\omterm{def:g2:water-on-earth:ocean}{océan} (environ 3\,700 mètres sous ce niveau)~;
une plage (environ 0 mètre)~; un village à 1\,500 mètres~; le sommet de
l'Everest (8\,849 mètres).
\end{solution}

\begin{solution}{exo:g4:world-physical-map:8}
Un fleuve est fait de nombreux ruisseaux qui se rejoignent, et l'on n'est
pas d'accord sur celui qui est son vrai commencement. Mesuré depuis une
source plus lointaine, un fleuve est plus long. Le Nil et l'Amazone ont des
longueurs si proches (environ 6\,700 à 6\,800 kilomètres chacun) que la
réponse dépend de ce choix.
\end{solution}

\begin{solution}{exo:g4:world-physical-map:9}
Le quart de 30 millions fait sept millions et demi de kilomètres carrés
(30 divisé par 4)~; le tiers fait 10 millions (30 divisé par 3). Le
Sahara, huit millions et demi, fait plus du quart de l'Afrique mais moins
du tiers.
\end{solution}

\begin{solution}{exo:g4:world-physical-map:10}
$2500 - 1000 = 1500$~: le refuge est 1\,500 mètres plus haut. Le refuge,
entre 2\,000 et 4\,000 mètres, est brun~; le \omterm{def:g1:landscapes-around-us:water}{lac}, à 1\,000 mètres, est à la
limite entre le jaune pâle (de 500 à 1\,000) et le brun clair (de 1\,000 à
2\,000)~: les deux réponses sont acceptées.
\end{solution}

\begin{solution}{exo:g4:world-physical-map:11}
10 kilomètres font 10\,000 mètres. $10000 - 8849 = 1151$~: l'avion est
plus haut que l'Everest, de 1\,151 mètres. Au-dessus de l'abîme
Challenger, il est à 10\,000 mètres au-dessus de la \omterm{def:g2:water-on-earth:ocean}{mer}, et le fond de la
fosse est à 11\,000 mètres sous la \omterm{def:g2:water-on-earth:ocean}{mer}~: $10000 + 11000 = 21000$ mètres
au-dessus du fond de la fosse.
\end{solution}
